\documentclass[UTF8,a4paper,11pt]{ctexart}
\usepackage[margin=2.2cm]{geometry}
\usepackage{graphicx}
\usepackage{booktabs}
\usepackage{enumitem}
\usepackage{titlesec}
\usepackage{hyperref}
\hypersetup{colorlinks=false,pdfborder={0 0 0}}

% 紧凑的章节标题
\titleformat{\section}{\large\bfseries}{\thesection.}{0.5em}{}
\titleformat{\subsection}{\normalsize\bfseries}{}{0em}{}
\titlespacing*{\section}{0pt}{1.0ex}{0.4ex}
\titlespacing*{\subsection}{0pt}{0.6ex}{0.2ex}
\setlength{\parskip}{0.3ex}
\setlength{\parindent}{2em}

\title{基于《小兵传奇》的星际战争游戏 Demo 设计}
\author{小组成员：高子豪（学号：2026011117）\quad 张容博（学号：2026011111）}
\date{}

\begin{document}
\maketitle
\vspace*{-2ex}

\begin{abstract}
\noindent 本项目拟以网络小说《小兵传奇》为世界观蓝本，使用 Unity 引擎与 Blender 资产制作，
开发一款以星际战争为主题、以阵营成长为主线的可玩游戏 Demo。设计的核心思路是让联邦、银鹰帝国、
家族三个可玩阵营共享同一套底层业务逻辑（关系网、作战体系与星际地图），
仅在初始条件与成长资源上产生差异，从而在有限工期内最大化内容复用。
本文说明选题依据、与已有作品的比较、系统方案与实现计划，并汇报已完成的设定分析、
系统设计、可点击原型与可编译的 C\# 工程骨架。
\end{abstract}

\section{选题与目标}

选择星际战争题材，一方面借鉴《戴森球计划》在生产与规模扩张上的正反馈体验，
另一方面利用《小兵传奇》已具备的完整设定降低世界观构建成本：小说给出了四种政体、
六等二十一级军衔、A--J 级舰船数值表、两大强国与无乱星系的割据格局，
以及从新兵到统一宇宙的四阶段主线。这些设定可直接转化为游戏的数值基线与剧情骨架。

我们计划交付一个可操作的游戏 Demo，支持三阵营中的至少一个完整可玩流程，
玩家通过军功（联邦、银鹰）或资源扩张（家族）逐步成长，并在后期进入相同的资源与政治博弈层。
AI 的应用落在游戏内部——以小说中的人格化中央电脑「星零」为原型设计 AI 副官，
提供战术建议与情报分析；本项目不以 AI 算法 novelty 为目标，评估标准是 Demo 本身的完成度与可玩性。

\section{相关工作与比较}

\begin{table}[h]
\centering
\small
\begin{tabular}{l l l}
\toprule
作品 & 核心机制 & 与本项目的差异 \\
\midrule
戴森球计划 & 自动化生产链、规模扩张 & 本项目以阵营博弈与叙事分支为核心，生产仅为后勤一环 \\
Homeworld & 舰队战术、立体空间指挥 & 本项目叠加关系网与成长系统，作战只是其中一层 \\
Stellaris & 4X 大战略、外交与派系 & 本项目聚焦单局角色成长，不做全局 4X 模拟 \\
\bottomrule
\end{tabular}
\end{table}

三者分别覆盖了生产、战术与战略，但都缺少「角色在派系政治中成长」这一叙事驱动层。
本项目的差异点在于：以小说剧情为剧本，将军功、资源与关系网数值绑定到剧情分支，
并让三个阵营复用公共模块以降低实现成本。

\section{方案设计}
主要的剧情线目前的设计中分成三条：联邦线、银鹰线、家族线。
\subsection{三阵营的差异与平衡}
联邦线以唐龙视角展开，前期由高层承担资源调度，玩家专注作战，需应对军功被冒领与派系倾轧；
银鹰线前期类似，但下属绝对忠诚，矛盾来自皇子夺嫡与贵族站队；
家族线从零开始经营资源与工业，自由度高但关系简单。
三条线前期水平相当：作战线省心于后勤，家族线省心于政治；
成长路线同构，后期玩家晋升高层后统一进入资源与政治博弈，系统合流。

\subsection{故事走向与资源调度}
主线按小说四阶段矛盾组织：新兵期、无乱星系立国、反唐联盟、统一宇宙。
设置 2--3 个改变结局的关键抉择，由关系网、军功与资源存量共同触发。
资源采用统一底层（资金、稀有金属、舰船、人口、能源），仅暴露时机不同：

\section{实现计划与分工}

\begin{table}[h]
\centering
\small
\begin{tabular}{l l l}
\toprule
阶段 & 任务 & 产出 \\
\midrule
第 1 周 & 纵向切片：真实战斗与 Unity 外壳 & 可操舵、开火、命中、结算的灰盒战斗 \\
第 2 周 & 战略最小闭环 & 配置化舰船属性、任务日志、两节点星图、存档 \\
第 3--6 周 & 公共模块与内容 & 关系网、经济、剧情的可运行版本 \\
第 7--10 周 & 阵营内容与副官 & 一个阵营的完整流程、AI 副官原型 \\
第 11--12 周 & 联调、测试与展示 & 可玩 Demo、演示视频 \\
\bottomrule
\end{tabular}
\end{table}

\noindent\textbf{分工}：成员 A（高子豪）负责真实战斗仿真，只改 \texttt{com.starpawn.combat} 与 \texttt{content/ships}；
成员 B（张容博）负责 Unity 可玩外壳、战略规则、关系网、任务剧情、AI 副官与最终装配，改其余包与 \texttt{Assets/Game}。
两人只在 \texttt{BattleStartSpec} 与 \texttt{SettleBattle} 交界，\texttt{CombatContracts} 暂时冻结；
第一周围绕同一场战斗并行开发，第五天在战斗结果接口处汇合。

\textbf{资源估计}：人力 2 人，工期约 12 周；技术栈为 Unity（运行时）+ Blender（资产制作，导出 glTF/FBX）。
主要资源需求为开发机器与美术素材，无额外算力或数据采购需求。

\textbf{前期成果}：已完成《小兵传奇》全文梳理与三份设定分析文档，制作了概念海报、系统架构设计，
以及含星图、关系网、作战、资源、剧情五个视图的可点击原型（作战页为实时 3D 战斗 HUD）。
工程侧按 12 个 Unity 包、21 个程序集搭起契约与假实现分离的骨架，编译期禁止跨模块引用，
21 个程序集已实际编译通过；14 个场景块的契约测试覆盖重复结算、失败回滚、战损继承、
快照隔离与内容校验。舰船、阵营、星图、任务、剧情与任务日志均已落为 JSON 并带引用校验，
改配置即可改变奖励与舰体属性。代码仓库见 \url{https://github.com/gaozihao1216/NLP-project}。

\section{AI 使用声明}
本 proposal 的撰写过程中使用了 AI 助手：用于梳理小说文本、提取设定并整理为分析文档，
以及辅助文字表述与结构组织。选题方向、系统设计决策与最终内容由小组成员确定并核对。

\end{document}
